\uuid{CTmr}
\titre{Matrice/Inversibilité/Rang}
\auteur{hanani ; mzari}
\organisation{exo7}
\licence{CC BY-NC-SA 4.0}
\datecreate{2020-01-20}
\niveau{L1}
\module{Algèbre}
\chapitre{Matrice}
\sousChapitre{Inverse, méthode de Gauss}
\difficulte{3}
\theme{Inverse d'une matrice, Matrice, Inversibilité, Rang}
\exo7qcmid{lille-2.matrices.033}

\contenu{
\begin{question}
\qclassification{108.05}
\qtags{motcle=Matrice/Inversibilité/Rang}
Soit $ M_n(\Rr)$ l'ensemble des matrices carrées d'ordre $n$  à coefficients réels  et $I$ la matrice identité.
Soit $A \in M_n(\Rr)$ telle qu'il existe un entier  $m \ge 1$ vérifiant :  $A^m+A^{m-1}+ \dots + A + I = 0$. Quelles sont les assertions vraies ?
\begin{answers}
\good{$A$ est inversible et $A^{-1} = A^m$}.
\good{$A$ est inversible et $A^{-1} = -(A^{m-1}+ \dots + A+I)$}.
\good{Le rang de $A$ est $n$}.
\bad{$A$ n'est pas inversible}.
\end{answers}
\begin{explanations} De l'égalité : $A \times  (A^{m-1}+A^{m-2}+ \dots + I)= -I$, on déduit que  $A$ est inversible et que $A^{-1}=-(A^{m-1}+ \dots + A+I)=A^m$. Puisque $A$ est inversible, le rang de $A$ est $n$.
\end{explanations}
\end{question}
}
